\color{blue}
\section{Proof of \Cref{prop:necessary-cond} \label{appendix:necessary}}

In our discussion of Gaussian channels in \Cref{sec:information}, we recalled that when an individual coordinate $X$ of unit variance suffers interference of power $P,$ it can convey at most $\ln(1 + P^{-1})/2$ nats to a receiver. In terms of mutual information, this means that
$$
I(X; X + Z) \le \frac 1 2 \ln (1 + P^{-1}) = \frac{1}{2 P} + O(P^{-2})
$$
for large $P.$ We begin by establishing a similar bound on mutual information between a superposition of Rademacher codewords and one of its summands.

\begin{proposition} \label{prop:mutual}
Let $F_1, \dots, F_k$ be $k$ independent Rademacher codewords, as in the definition of a Rademacher dictionary. Then for any random variable $Y$ independent from $(F_1, \dots, F_k)$ and supported on the set $\{-1/\sqrt{d}, 1/\sqrt{d}\}^d,$
$$
I\left(Y; Y + \sum_{i = 1}^k F_i \right) \le d \left(\frac 1 {2 k} + O(k^{-2})\right),
$$
where the remainder of $O(k^{-2})$ depends only on $k.$
\end{proposition}

\begin{proof}
Let $Z = F_1 + \dots + F_k.$ By independence,
$$
I(Y; Y + Z) = H(Y + Z) - H(Y + Z | Y) = H(Y + Z) - H(Z).
$$
Now, let $\pi_i$ be the projection onto the $i$th coordinate. Since the coordinates of $Z$ are independent,
$$
H(Z) = \sum_{i = 1}^d H(\pi_i(Z)),
$$
Furthermore, by subadditivity,
$$
H(Y + Z) \le \sum_{i = 1}^d H(\pi_i(Y + Z)).
$$
The entropy of $\pi_i(Y + Z) = \pi_i(Y) + \pi_i(Z)$ is unaffected by applying a rescaling where $\pi_i(Y) \in \{0, 1\}$ and $\pi_i(Z)$ has a binomial distribution with $k$ trials. Overall, it is enough to prove that
$$
I(X; B_k + X) = H(B_k + X) - H(B_k) \le \frac{1}{2 k} + O(k^{-2})
$$
where $B_k$ has a Binomial distribution with $k$ trials and $X$ is independent from $B_k$ and supported on $\{0, 1\}.$ Furthermore, it is easy to check that $H(B_k + X)$ is maximized where $X$ is uniformly distributed over $\{0, 1\},$ so it suffices to prove our inequality in this case.

The posterior on $X$ conditional on $B_k + X$ is given by
$$
\Pp(X = 1|B_k + X) = \frac{B_k + X}{k + 1}.
$$
Putting $S = (B_k + X)/(k + 1) - 1/2,$ we find that
$$
I(X; B_k + X) = H(X) - H(X|B_k + X) = \ln 2 - \E \left[ h\!\left(S + \frac{1}{2}\right)\right].
$$
Taking a series expansion at $1/2$ gives
$$
I(X; B_k + X) = \E \left[ 2 S^2 + O(S^4)\right].
$$
Odd moments of $S$ vanish, and the first even moments are
$$
\E[S^2] = \frac{1}{4(k+1)}, \quad \E[S^4] = \frac{3k+1}{16(k+1)^3}.
$$
We conclude that
$$
I(X; B_k + X) = 2 \E[S^2] + O(\E[S^4]) = \frac{1}{2(k+1)} + O(k^{-2}) = \frac{1}{2k} + O(k^{-2}).
$$
\end{proof}

It will also be useful to state Fano's lemma in the following way.

\begin{lemma} \label{lemma:fano}
Let $X$ be uniformly distributed over $[N]$. Then for any random variable $Y$ and any function $f,$
$$
\Pp(f(Y) \neq X) \ge 1 - \frac{\ln 2 + I(X | Y)}{\ln N}.
$$
\end{lemma}

\begin{proof}
A usual statement of Fano's lemma is that, if there exists some function $f$ satisfying $f(Y) = X$ with probability $p,$ then
$$
H(X|Y) \le h(p) + (1 - p) \ln N,
$$
where $h(p)$ is the binary entropy function. Since $h(p) \le \ln 2$ and $H(X) \le \ln N,$ this implies that
$$
p \le 1 - \frac{H(X|Y) - \ln 2}{\ln N} = \frac{\ln 2 + \ln N - H(X|Y)}{\ln N} \le \frac{\ln 2 + I(X; Y)}{\ln N}.
$$
which implies our statement.
\end{proof}

Now we proceed to the proof of \Cref{prop:necessary-cond}.

\color{blue} \begin{proposition*}
    Let $C > 0$ be arbitrary. Over any regime where
    $$
    d \le C k \ln N, \quad \omega(1) \le k < N/2,
    $$
    for Rademacher dictionaries it holds that
    $$
    \liminf_{N \to \infty} b(d, k, N; \text{top-$k$}) \ge 1 - \frac C 2.
    $$
\end{proposition*}

\color{blue} \begin{proof}
We claim that for all integers $d, k, N \ge 1,$
$$
b(d, k + 1, N + k; \text{top-}k) \ge 1 - \frac{\ln 2}{\ln N} - \frac{d}{2 k \ln N}(1 + O(k^{-1})).
$$
Indeed, substituting variables will then give
$$
b(d, k, N; \text{top-}k) \ge 1 - \frac{\ln 2}{\ln (N - k + 1)} - \frac{d}{2 (k - 1) \ln (N - k + 1)}(1 + O(k^{-1})).
$$
Under the regime $d \le C k \ln N$ and $\omega(1) \le k \le N/2,$
$$
\frac{d}{2(k - 1) \ln(N - k + 1)}(1 + O(k^{-1})) = \frac{d(1 + O(k^{-1}))}{2 k \ln N (1 + o(1))}\le \frac{C}{2} + o(1).
$$
We conclude overall that
$$
b(d, k, N, \text{top-}k) \ge 1 - o(1) - \frac{C}{2},
$$
which proves our result.

We now prove the claim. We define three independent random variables:
\begin{itemize}
\item $F \in \R^{d \times (N + k)}$ is a Rademacher dictionary;
\item $K$ is uniformly random in $[N]$;
\item $\sigma$ is a random permutation of $[N + k]$.
\end{itemize}
Define the vector
$$
X_i = \begin{cases}
1 : i = K \lor i > N \\
0 : \text{otherwise}.
\end{cases}
$$
We write $(F^0, F^1)$ for the first $N$ and last $k$ coordinates of $F,$ and similarly write $(X^0, X^1)$ for the first $N$ and last $k$ coordinates of $X.$ Thus, the superposition code for $X$ under $F$ can be written as
$$
Y = F X = F^0 X^0 + F^1 X^1.
$$
We will write $\sigma(X)$ and $\sigma(F)$ for the permutations of $X$ and $F$ under $\sigma$ defined in the natural way. In these conditions, $(\sigma(X), \sigma(F))$ is an independent pair of a uniformly random $k + 1$-sparse vector and a Rademacher dictionary.

Now, take a function $g \colon \R^{d \times N} \times \R^d \to [N]$ so that $g(A, y)$ is any coordinate where the maximum value of the vector $y A$ is attained. If a top-$k$ decoder succeeds at decoding $\sigma(X)$ from $Y$ using the dictionary $\sigma(F),$ then $F X^0$ satisfies
$$
\langle F^0 X, Y \rangle \ge \langle F_i, Y \rangle
$$
for all $i \in [N].$ It follows that $g(F^0, Y) = K,$ except potentially when multiple coordinates of $Y F^0$ attain the maximum value. However, in this situation, the probability of the success of top-$k$ is no greater than the probability that $g(F^0, X) = K.$ Therefore,
$$
b(d, k + 1, N + k; \text{top-}k) \ge \Pp(g(F^0, F X) \neq K).
$$
Applying \Cref{lemma:fano} to the map $g,$ have
$$
\Pp(g(F^0, FX) \neq K) \ge 1 - \frac{\ln 2}{\ln N} - \frac{I(K; (F^0, FX))}{\ln N}.
$$
Finally, since
$$
I(K; (F^0, FX)) = I(K; FX \mid F^0) \le I(F^0 X^0; FX \mid F^0)
$$
and since $FX = F^0 X^0 + F^1 X^1$ is a linear combination of the codeword $F^0 X^0$ with $k$ independent Rademacher vectors, by \Cref{prop:mutual} we have
$$
I(K; (F^0, FX)) \le d \left(\frac 1 {2 k} + O(k^{-1})\right).
$$
Altogether,
$$
b(d, k + 1, N + k; \text{top-}k) \ge \Pp(g(F^0, FX) \neq K) \ge 1 - \frac{\ln 2}{\ln N} - \frac{d}{2 k \ln N}\left(1 + O(k^{-1})\right)
$$
which proves our claim.
\end{proof}

\color{black}
